\section{Exegetical Alternatives Tested}

\subsection{Only absence, no positive sign}

One reading treats the cloths as inventory: they prove only that the body is
gone. This properly resists imaginative geometry, but it underreads John's
slowness. The Evangelist distinguishes the two cloth groups, their relation,
three acts of seeing, entrance order, and belief. Details can remain
multivalent without being narratively inert.

\subsection{A deliberately tidied tomb}

Another reading imagines the risen Jesus folding or rolling the face-cloth as
an act of calm. It captures the contrast with hurried theft, but John never
names the agent who arranged anything after burial. The participle can describe
the cloth's wrapped condition from burial itself. Calm is a reasonable
inference from differentiated order, not a narrated gesture.

\subsection{Evidence designed as a message}

Chrysostom approaches this: the disposition is providentially intelligible and
moves disciples toward faith. A stronger modern version makes every cloth a
deliberate encoded communication. That exceeds the text. A sign can be
providential without operating as a cultural cipher. John's own explanation is
the transition from seeing to believing, not an appended legend about how to
decode folds.

\subsection{Proof of a glorified body's mode}

The collapsed-wrappings account offers a physical mechanism: Jesus rose through
the material. Catholic theology affirms a glorified body no longer subject to
ordinary constraints, but John 20 does not describe the instant of
Resurrection. The later closed-door appearances provide an analogue, not a
camera angle into the tomb. Doctrine establishes possibility; it does not turn
a possible reconstruction into exegesis.

\subsection{A relic-identification text}

John's nouns naturally interest readers of relic traditions. This study,
however, has not established provenance from the tomb to any surviving cloth.
Vocabulary comparison can test whether a proposed object is compatible with a
text; compatibility cannot establish identity. Historical custody, material
dating, contamination, conservation, and scientific disputes belong to a
separate work.

